\section{Introduction}

% An emerging line of research has begun to explore both fronts of this question. 
% On the methodological side, harness evolution has progressed from prompt and workflow search to scaffold code refactoring, and further toward the end-to-end optimization of executable harnesses.
% On the evaluation side, existing work focuses on the assessment of harness evolution methods~\cite{seaeval,seagym} and autonomous agent development~\cite{mac}, as well as in-depth analyses of the underlying causes of harness benefits~\cite{seaeval,seagym,agentenvcoevolution}.
% Collectively, these parallel lines of research demonstrate how harness evolution can be achieved and characterize what successful evolution entails.
% However, systematically benchmarking an LLM's intrinsic capacity to improve harnesses introduces unique obstacles that existing studies leave unaddressed.

% \textcolor{blue}{
% An emerging line of research has begun to explore both fronts of this question. 
% On the methodological side, harness-evolution techniques attempt to search over prompts and workflows, refactor scaffold code, or optimize executable harnesses end-to-end~\cite{adas,aflow,stop,godelagent,dgm,metaharness,harnessx}.
% On the evaluation side, studies on self-evolving agents evaluate self-evolutionary dynamics, capability transfer, and long-term adaptation~\cite{seaeval,harnessupdating,evoagentbench,agentenvcoevolution,seagym}. 
% Collectively, these parallel lines of research demonstrate how harness evolution can be achieved and characterize what successful evolution entails.
% However, systematically benchmarking an LLM's intrinsic capacity to improve harnesses introduces unique obstacles that existing studies leave unaddressed.
% }

% Specifically, benchmarking this capability requires addressing three key challenges.

% localized prompt tuning.


% \textcolor{blue}{
% Specifically, reliably benchmarking harness-evolving capability requires addressing three key challenges.
% First, \emph{harness sensitivity}: benchmark tasks should be responsive to harness improvements rather than dominated by base model capability, as gains on harness-insensitive tasks provide weak evidence of genuine harness evolution.
% Second, \emph{cross-split generalization}: validation and evaluation tasks should exhibit aligned responsiveness to harness variation. Naively selecting tasks may lead to overfitting to a specific task set rather than transferable harness improvement~\citep{rethinkingharnesseval}.
% Third, \emph{long-horizon evolution}: harness evolution is a long-horizon iterative process that requires sustained refinement across multiple rounds, where models diagnose failures, formulate improvement hypotheses, and progressively update executable harnesses.
% Without addressing these challenges, benchmark gains may conflate harness-evolving capability with model strength or task-specific optimization.
% }


Recent breakthroughs in large language models (LLMs) have remarkably elevated the capacity of agentic systems to execute complex, long-horizon tasks at an end-to-end, system-level scale~\cite{deepseekv4,glm5,anthropic2026opus,anthropic2026fable}. These gains arise not only from improvements in the underlying models, but also from carefully engineered agent harnesses that structure, coordinate, and constrain agent behavior, as exemplified by Claude Code~\cite{claudecode} and Codex~\cite{codex}.
In this context, an increasing body of work therefore seeks to explore  the potential for increasingly capable agentic systems to achieve self-improvement, including autonomous scientific research~\cite{alphaevolve}, automated model post-training~\cite{posttrainbench}, and autonomous harness evolution~\cite{dgm,metaharness,harnessx}. Among these lines of research, harness evolution is generally considered the first step toward achieving self-improvement~\cite{weng2026harness}.
Consequently, a pivotal question emerges for the foundational models: \emph{can language models truly improve agent harnesses, and how can we systematically benchmark this evolutionary capability?}

Recent work has explored harness evolution on two key fronts: methodologically, progressing from prompt search to end-to-end executable scaffold refactoring; and evaluatively, benchmarking evolution methods~\cite{seaeval,seagym} and agent development~\cite{mac} while analyzing underlying harness benefits~\cite{seaeval,seagym}. Collectively, these studies demonstrate how harnesses evolve and what constitutes success. However, systematically benchmarking an LLM's intrinsic capacity for harness evolution introduces unique, unaddressed obstacles.

Specifically, we identify three key challenges.
First, \emph{harness sensitivity}: benchmark task performance must be responsive to harness improvements rather than dominated by base model strength.
Second, \emph{cross-split generalization}: validation and evaluation splits must exhibit aligned responsiveness to prevent task-specific overfitting~\citep{rethinkingharnesseval}.
Third, \emph{long-horizon evolution}: models must sustain multi-round iterative refinement---diagnosing failures, formulating hypotheses, and progressively updating code.
Unaddressed, these challenges conflate true harness-evolving capability with model strength or task-specific overfitting.

To address these challenges, we introduce \textbf{Evo-Bench}, the first benchmark evaluating models' intrinsic harness-evolving capability. Evo-Bench adopts a controlled, long-horizon research setting with a fixed policy model across three domains (Search, Office, and General Agent tasks) over five established benchmarks~\cite{browsecomp,hle,gdpval,apexagents,claweval}. To rigorously construct this benchmark, we propose a novel \emph{harness-guided benchmark construction framework}. Built upon this framework, Evo-Bench incorporates three core designs: (1) \emph{harness-sensitive task construction}, identifying tasks responsive to framework improvements; (2) \emph{sensitivity-aware stratified splitting}, ensuring strict distributional consistency and disjoint validation and evaluation suites; and (3) \emph{controlled long-horizon evolution}, enabling sustained optimization of agent harnesses.
% multi-turn, evidence-driven iterative harness evolution.

In our main evaluation across 9 frontier and open-weight models on \textbf{Evo-Bench}, top evolvers like GPT-5.6 Sol and Claude Opus 4.8 achieve substantial gains of 16.6 and 16.1 over the seed harness, respectively, closely approaching the human-engineered baseline of 47.5. Our analysis yields three key insights:
(1) \emph{evolutionary behavior exhibits early saturation}, with models rapidly discovering high-quality structures before introducing detrimental modifications in later rounds;
(2) \emph{gains are highly domain-dependent}, as evolvers effectively replicate web navigation in Search for massive gains and can even surpass manual engineering in General tasks, but struggle in Office tasks that demand highly specific processing workflows; and
(3) \emph{synthesized harnesses act as transferable reasoning structures}, demonstrating cross-policy robustness by driving consistent gains across diverse policy models from Qwen, DeepSeek, and GLM.

To summarize, our contributions are three-fold:

\begin{itemize}
\item We introduce \textbf{Evo-Bench}, the first benchmark evaluating LLMs' intrinsic harness-evolving capability---their ability to autonomously refine executable code harnesses across diverse domains beyond static task solving.

\item We propose a harness-guided benchmark construction framework that first induces diverse harnesses through auxiliary-task evolution and then constructs the final benchmark via harness-guided task selection.

\item We provide a systematic scientific account of harness evolution across 9 frontier models, characterizing substantial gains reaching 16.6 points, temporal anomalies like early saturation, and the robust transferability of evolved structures across distinct policy architectures.
\end{itemize}

% \item \textcolor{blue}{We introduce Evo-Bench, the first benchmark for comparing harness-evolving capability across agentic models, defined as the ability to autonomously evolve general-purpose harnesses through long-horizon engineering across diverse tasks, benchmarks, and domains, beyond conventional task-solving evaluation.}


% \item \textcolor{blue}{We conduct a comprehensive evaluation of 9 frontier models, quantifying substantial optimization gains (up to +16.6), uncovering critical behavioral dynamics like early saturation, and demonstrating robust harness generalization across diverse target policy architectures.}
